\section{Which countries matter in the public proxy?}
\label{app:country-combinations}

This appendix asks whether particular country pairs or triples dominate the public delivered-easing association. The public policy-rate exercise does not reveal which curves created the baseline fifteen inversion episodes. The analysis uses the same BIS policy-rate and exchange-rate data as Appendix~\ref{app:public-checks}.

\subsection{Combination screen}

All 36 currency pairs and 84 triples are enumerated. For a given combination, an onset occurs when every member cuts its policy rate by at least 10 basis points after three months without that same joint cut. The outcome is the public shadow-carry spot return accumulated over event months zero and one. The ranking of all nine currencies uses lagged policy-rate differentials; the tested country combination defines the event, not the portfolio.

Reference values are reported only for combinations with at least six valid events. Every eligible event sequence is shifted through every circular calendar rotation. Raw two-sided rotation references are supplemented with a common-rotation maximum absolute standardized reference across every eligible pair and triple. Under the maintained null that simultaneous cyclic shifts are exchangeable, the maximum reference controls the declared family; otherwise it is a finite family diagnostic. The 120-combination family was declared in code and exhaustively reported but was not preregistered. Episode-resampling intervals, leave-one-event means, and alternative minimum-event rules diagnose small-sample fragility.

\begin{table}[H]
\centering
\caption{Country combinations in the public delivered-easing proxy}
\label{tab:public-country-combinations}
\small
\begin{tabular}{lrrrr}
\toprule
Combination & Events & Mean $[0,1]$ & Raw ref. & Max-$|z|$ ref. \\
\midrule
CHF--GBP & 6 & -4.54 & 0.009 & 0.035 \\
AUD--EUR & 6 & -4.26 & 0.004 & 0.055 \\
GBP--SEK & 9 & -3.32 & 0.013 & 0.119 \\
EUR--NZD & 6 & -2.80 & 0.053 & 0.486 \\
CHF--SEK & 7 & -2.74 & 0.070 & 0.508 \\
CAD--NOK--NZD & 7 & -4.09 & 0.013 & 0.066 \\
GBP--NZD--SEK & 8 & -3.71 & 0.018 & 0.086 \\
NOK--NZD--SEK & 6 & -3.62 & 0.026 & 0.235 \\
CAD--GBP--NOK & 6 & -3.33 & 0.031 & 0.321 \\
AUD--NOK--NZD & 6 & -2.87 & 0.048 & 0.519 \\
\bottomrule
\end{tabular}
\begin{minipage}{0.94\textwidth}\footnotesize
\textit{Notes:} The table reports the five most negative eligible pairs and triples. An event requires every listed country to cut its BIS policy rate by at least 10 basis points after three months without the same joint cut. The outcome is the cumulative policy-ranked log spot-return proxy in event months 0 and 1, in percentage points. Portfolio legs target three currencies per side and expand to include all cutoff ties with equal weight. Combinations require at least six valid events. Raw reference values use every circular timing rotation. The final column is the common-rotation maximum absolute standardized reference across every eligible pair and triple. Under simultaneous cyclic-shift exchangeability it controls the declared family; otherwise it is a finite family diagnostic. The family was not preregistered. This is a delivered-easing proxy, not a replication of the yield-curve state.
\end{minipage}
\end{table}


Of 120 combinations, 42 have at least six events. CHF--GBP is the only combination meeting the 5-percent common-rotation criterion: six joint-cut onsets have a mean cumulative two-month shadow-carry log spot-return proxy of $-4.54$ percentage points, raw reference $0.009$, and maximum-$|z|$ reference $0.035$. The 90-percent episode-bootstrap interval is $[-9.52,0.31]$, however, and deleting October 2008 reduces the mean to $-2.52$. AUD--EUR is the next-most-negative pair, at $-4.26$ with maximum-$|z|$ reference $0.055$. CAD--NOK--NZD is the most-negative triple, at $-4.09$ with reference $0.066$.

The family result depends on admitting sparse combinations. With a six-event minimum, 42 combinations are eligible and one meets the 5-percent reference criterion. With minimums of eight or ten events, 25 and 15 combinations remain, no result meets that criterion, and the smallest maximum-$|z|$ reference values are 0.064 and 0.229. The evidence does not establish a stable country bloc.

\subsection{Sensor countries versus loss-bearing currencies}

The public lead does not match the baseline loss-bearing cross-section. CHF and GBP have near-zero or modest equity betas and small, imprecise state-conditioned currency losses. AUD and NZD bear the larger losses. The own-curve horse race also shows that a country's own curve generally does not forecast its currency after conditioning on the synchronized state.

This contrast suggests a testable distinction: countries whose bond or policy states reveal a shared episode need not be the currencies that bear its losses. Joint Swiss--U.K. easing may function as a calendar sensor for broad stress, while high-rate, high-beta currencies absorb the carry loss. The distinction is economically consistent with synchronization and exposure ordering, but the present result is only hypothesis-generating because delivered easing is downstream of forward-looking yield-curve inversions.

The direct test requires the nine-country live-state matrix. A fifteen-episode-by-nine-country contributor ledger could enumerate which curves activate each episode, then compare contributor combinations with the subsequent currency-loss cross-section under family-wide inference. The public combination screen specifies that future test but does not decompose the headline result.
